
% 来源：EPUB part109.xhtml
\section{用户导向}

很久以前有个记者朋友来我这问我一个问题：“为什么我采访了好几家企业，他们都说自己是电话机行业里的第二？到底谁第一呢？”我说，我们知道但不说，因为消费者不关心这个问题。（2010-06-08）

% 来源：EPUB part110.xhtml
\subsection{自己不喜欢的东西别拿出来卖}

\BookQuestion{网友：以消费者为导向，您是怎么想到这个产品哲学的，特别想了解其中的逻辑？}

消费者导向这个想法很自然的。设计产品不是简单希望用户买，而是要想到人家买回去干什么。有些人的生意不太好或者没办法长期好的原因往往是太在意眼前生意而忽略了用户的真实需求。没有哪个企业是总能找对用户的需求点的，所以再好的企业也会有失败的产品。但那些一直努力聚焦在用户需求的公司，时间长了往往会比那些生意导向的公司表现好。（2020-10-21）

\BookQuestion{网友：对消费者导向有点困惑，有好书推荐吗？}

这东西还需要看书？最通俗的说法就是自己不喜欢的东西别拿出来卖。如果你觉得这句话对，你根本不需要看书。如果你觉得不对的话，其实看啥书都没用。

当然自己喜欢的东西未必就一定好卖，那是因为经常没办法让别人知道你喜欢的原因，或者是因为市场营销费用太高，或者就是因为产品教育太难（比如有些工程师喜欢的东西就没办法让普通老百姓喜欢）。（2012-04-28）

\BookQuestion{网友：您早年做企业的时候，怎么把消费者导向和结果导向深入人心，在做事上不变形的？}

事实上我们也会经常有点小变形的，发现了就改。利益导向就像打球时的加力，是个心魔，克服不容易。好的企业文化对克服这个心魔作用很大。（2020-10-12）

\BookQuestion{网友：您怎么看企业文化中的用户导向、结果导向和追求极致？三者之间有什么逻辑关系？}

用户导向是一种利润之上的追求，简单讲就是想着用户到底要的是什么，我们如何去实现。结果导向指的是要把事情做对。追求极致还是要把事情做对的范畴。

对一个公司而言，去找到（开发）能够长期满足用户需求的产品就是做正确的事情。如何找到是如何把事情做对的范畴。（2019-09-26）

\BookQuestion{网友：我们是美国亚马逊上的一个中国商家，销售服装。您觉得这个生意有多大概率能做10年甚至更久？}

能做多久应该取决于你自己。如果你一直想着你的用户的需求，10年或更久为什么会是问题？难道10年后大家就不穿衣服了？很久以前黄峥曾经问过我一个问题：用户导向是放之四海而皆准的吗？我说，我相信是这样！我想他懂了！（2019-05-15）

\BookQuestion{网友：我有缘在OPPO工作了五年。这五年，让我对产品、品牌、营销和渠道都有了一定的理解。请问OPPO的这一套打法不适用于哪些行业或产品？}

消费者导向放之四海而皆准。技巧的东西谁都可以学，但我们之所以成为我们，除了我们一直努力做的事情外，我们不做某些事情的决定也同样重要。（2019-06-10）

\BookQuestion{网友：企业以消费者为导向是做消费者喜欢的产品，以市场为导向是做市场喜欢的产品。市场是谁？能否说说市场导向与消费者导向的区别是什么？}

从字面上理解，消费者导向指的是企业在做产品决策的时候是基于产品最后带给消费者的体验的考虑。市场导向看上去是基于现在是否好卖的。

很多时候这两者区别不是那么明显，因为好卖的东西往往就是消费者喜欢的。但消费者往往只是根据现有的产品来决定自己的喜好，所以经常会显得短视，如果企业总是市场导向的话，则陷于短视状况的概率就会相对大些。而总是注重消费者体验的公司则往往会眼光放得长远些，出伟大产品的概率就会大很多。

过去的国营企业往往基于计划经济，所以连市场导向都做不到。绝大部分企业最后大概只能到达市场导向的地步，因为大部分企业都是利润导向的，而要做到消费者导向需要有利润之上的追求，或者说需要能够放弃一些短期利益去做出最好的产品来满足消费者真实的需求。其实有些企业或多或少也是有利润之上的追求的，但能力上还没有达到做出伟大产品的地步，所以这类企业看起来可能没有那么强大，但生命力却往往会比人们想象的强很多。

只有极少数的企业是真的在骨子里具有消费者导向同时又具备实现消费者导向的能力的，这些公司最后往往会被人们称为伟大的公司。没有利润之上的追求的公司不大会真的有消费者导向的文化，最后也不大可能成为伟大的公司。具有利润之上追求的公司往往多少具有了消费者导向的文化或叫基因，在某些条件下就有可能会慢慢变成伟大的公司。

公司是由人来经营的，所以伟大的公司也不是一成不变的，时间有可能会改变一些东西，让一些伟大的公司变回到不那么伟大。（2013-04-07）

\BookQuestion{网友：利润导向从表面看是没什么错的，就好像一个家庭的男人要拼命给家里搞钱养活老婆孩子。如果不以利润为导向，男主人拼命把自己所从事的工作搞到最棒，其实钱的问题会自然解决掉。}

如果没有“取之有道”的约束，利润导向很容易导致企业不择手段。国内现在这个问题实际上非常严重。（2010-06-09）

\BookQuestion{网友：我特别在意企业是不是行业第一。}

其实是不是第一没有那么重要，因为消费者在买东西时一般并不在乎谁是第一，他们在乎的是买的东西是什么（消费者导向）。为追求市场占有率而追求市场占有率是大多数企业最容易犯的错误。（2010-06-07）

这里不在乎排名的意思是我们只在乎用户的感受，排名只是结果而已，不是追求的目标。（2010-06-09）

% 来源：EPUB part111.xhtml
\subsection{长期而言，消费者作为群体是理性的}

\BookQuestion{网友：您认为顾客是理性还是非理性的？您是满足顾客，还是让顾客忠诚于公司产品？}

长期而言，消费者作为一个群体是理性的。我们无法让消费者“忠诚于”我们公司的产品，但我们应该尽力去发现消费者的需求并尽量满足。所谓的忠诚度实际上应该是消费者长期的满意度的积累。所以长期满意度越高则忠诚度越高。要积累很高的满意度是需要很长很长的时间的，但破坏则只需要很短很短的时间。（2010-10-29）

\BookQuestion{网友：长期来说消费者会知道产品价值，但短期不认同的情况下企业应该怎样平衡呢？}

短期内不是所有的消费者认同，但总是有认同的。（2010-10-30）

\BookQuestion{网友：我认为中国很多消费者还不够理性，特别在一些需要专业知识的产品上。}

作为一个群体，消费者绝对是理性的。但这种理性不是每时每刻体现在每一个消费者身上，其会表现为不好的东西最后一定不会卖得好！不过，这并不会让好东西自然变得好卖。好酒也怕巷子深啊。（2010-03-25）

哪里的消费者都是一样的。我们认为消费者是理性的意思是从长期来看的，套用一句俗语叫“童叟无欺”。也就是说无论消费者眼前是否理性，我们都一定要认为他们是理性的。不然的话，你经营企业就可能会有投机行为，甚至会有不道德行为。（2010-03-26）

顾客作为整体绝对是理性的，只要时间足够长。没有这个理解的人，是很难理解生意的。（2023-12-23）

\BookQuestion{网友：今天安装了一个鲁班大师，在补丁更新的时候360竟然提示会导致系统崩溃！我就真不信！结果系统好好的速度还更快了。我依次卸了360杀毒、360安全卫士、360保镖。这下电脑清净了。提供服务的企业太霸道了，我是接受不了！}

能真的卸载就不错了。以前装过一个叫3721的东西，装完就发现不对，马上想卸载却无论如何找不到卸载的办法，最后只好重新装了一次机。腾讯要厚道些，我装过3次QQ，都是一装就发现被加了很多东西，然后我就卸载了，不过好像有一次也有点什么卸不干净的，但没认真查。以前装过一次迅雷，也是一装就被强行装上了很多东西，非常不舒服，不过卸载很方便。我也不太明白为什么我们这么多企业对自己那么没信心，非要做那么短视的事情。（2010-11-06）

\BookQuestion{网友：近两年我最喜欢的广告语是“充电五分钟，通话两小时”，简单明了，直指要点，朗朗上口。}

哈，有人说“充电五分钟，通话三个半小时”，然后他破产了，为什么呢？（2018-08-24）

% 来源：EPUB part112.xhtml
\subsection{站在用户的立场想问题}

\BookQuestion{网友：您做生意的时候，会不会和对手比较，您会选一个公司作为榜样或者追赶的对象吗？}

最重要的是用户的需求或叫体验，其他都是浮云。（2013-07-04）

\BookQuestion{网友：您公司的企业文化中有两处我不能领会：1. 通过科学、严谨的市场调查，充分研究消费者的需求，一切工作须以消费者的真实需求为原点而展开。（为何不直接写消费者需求，而写成消费者真实需求？）2. 在公司内部的日常工作中，要坚持内部客户导向的原则。（何为内部客户？）}

呵呵，玄乎地讲就是用心。用户的需求不一定是问出来的，经常是自己体会出来的。iPhone就是一个典型的例子。

1. 大概是为了把消费者需求和客户需求分开。客户需求有时候和消费者的真实需求并不完全重合。

2. 内部客户大概就是工作中要打交道的所有人。（2010-10-24）

\BookQuestion{网友：如何敏锐地发现用户需求，而且是真实的潜在的需求？}

发现用户需求的诀窍可能是总是站在用户的立场想问题。（2010-03-30）

\BookQuestion{网友：iPad上市之前，乔布斯和媒体记者在一起，被问到在开发新产品时苹果做了什么样的消费者和市场调查，这些调查起了什么作用。“我们没有做任何调查”，乔布斯回答说，“知道自己想要什么，不是消费者的工作。”如何理解乔布斯的话？}

他说得很对，因为他知道消费者想要什么。知道消费者想要什么是一件非常困难的事，没听说市场调研可以解决这个问题。市场调研只能够帮你印证你的想法而已，你如果自己没想法，市调是绝对没用的。其实苹果还是会做一些市调的，比如偶尔放点消息出来看看反应什么的。（2011-02-26）

\BookQuestion{网友：我一直不相信市调公司的数据，他们为钱做的数据在我看来一定会有偏差。我更相信自己的直觉。}

市调公司应该是中立的，是否有偏差取决于水平。你如果把问卷完全交给市调公司设计的话，得出的结论容易出偏差。我的观点是，市调结果是用来验证你的直觉的。如果你的直觉已经很清晰，其实用不着市调。但当你觉得没把握时，市调有时能帮很大忙。（2010-03-24）

\BookQuestion{网友：据称阿里某高管在内部强调，如果只能听一个人的声音那就是用户的声音，将告别过去以GMV（Gross Merchandise Volume，商品交易总额）导向的发展模式转向为以消费者为中心。大道对此有何看法吗？}

“如果只能听一个人的声音那就是用户的声音”，其实如果还能听第二个人的声音，那也应该是消费者的声音，依此类推。（2023-04-23）

“我们砸了自己的脚，我们过去忽视了用户体验，过去几年里，阿里落后了，淘天集团的问题在于忽视用户体验，整个集团已经为此进行了重组……”（阿里巴巴联合创始人蔡崇信专访）

也许“让天下没有难做的生意”要做一些调整了，因为这句话里面最终体现不出来用户？

从第一次看到这句就觉得不太合适，我很久以前应该在雪球也提过。

\BookQuestion{网友：可不可以理解为用户导向是所有生意的基础和本质，也是所有商业模式的核心，是在做正确的事情？}

最终用户是关键。就像花生是地里长出来的而不是商店里买的一样。（2024-04-04）

\BookQuestion{网友：狗粮的消费者是狗还是狗主人呢？中兴选的是狗主人，真正的消费者导向，应该是选狗，也有可能运营商需求包含了真正的消费者需求。}

经典案例啊。（2010-07-07）

\BookQuestion{网友补充：狗不买狗粮，但狗粮是狗吃的，狗才是真正的消费者。如果狗主人凭自己喜好买狗粮，而狗不喜欢吃，这种狗粮终究还是不好卖的。狗粮的口味（消费者）考察应该是选不同类型的很多狗试验，而不是仅仅在包装上迎合购买狗粮的狗主人。}

\BookQuestion{网友：你以前说过客户与用户有时是有区别的，但我到现在也搞不清头绪。}

有个狗粮的例子，就是客户满意但用户不满意的典型。（2012-03-03）

以客户为中心就是生意导向的意思，和以用户（消费者）导向的意思很不一样。（2024-01-14）

\BookQuestion{网友：能做到市场导向已经很不错了，可以做到消费者导向，并坚持下来更加是了不起！}

能坚持下来确实非常不容易，苹果也是几十年后才有今天的。（2013-04-11）

\BookQuestion{网友：我是中兴通讯手机事业部负责欧洲市场推广的。中兴手机这几年进步也很大，主要优势在于快速定制，满足运营商需求，抓住了类似国内中国移动、电信的大客户捆绑销售，走了跟步步高不同的路线。问题在于基本放弃了渠道市场，在某种程度上忽视了消费者需求，没有做到消费者导向，而是运营商导向。中兴手机的发展方向应该是学习中国台湾的HTC，从ODM（Original Design Manufacturer，原始设计制造商）逐步到自有品牌，从单一追求销量到追求销售附加值。道理似乎大家都懂，但做起来很难，是不是企业的文化和基因很难改变呢？每年为了追求销量忘记或者忽视了一些最根本的东西。也许是因为上市公司当期业绩压力大吧。}

上市公司当期业绩压力大的说法很有趣（尽管无数上市公司都这么说）。难道大家愿意让下期压力更大？（2010-07-07）

\BookQuestion{网友：国产手机的售后服务黑幕重重，各个品牌的售后基本都是外包给某个公司。这些公司为了利润最大化，完全没有道德底线可言，这对品牌的伤害会很大。难道厂家是完全不知情吗？}

别人不知道，我们的服务绝对是我们自己的。（2013-03-17）

\BookQuestion{网友：最近几年多次找格力的售后维修空调，需要更换的零件比如电路板、压缩机等都没有现成的，要向公司申请调货，有时候两三天能到，有时候要等一个星期。我搞不懂为什么这个环节一直没改进。空调坏的时候一般是大热天，维修速度对用户来说非常重要。}

没做过这种大件，不好评价太多。当年曾经想过空调的生意，就是想不通怎么才能做好安装和售后才放弃的。我当年提的售后服务条款里有一条，叫“立等可取”。当时修不好的可以临时给一台先用着，大件这么做可能有点困难？以前我在这里说过一个例子，我去苹果店，说买新机没有，但坏了可以换个新的。空调似乎更应该如此啊，大热天要命的。（2013-09-06）

\BookQuestion{网友：品类很少才有可能做到苹果那样吧。你们公司的产品品类不少，售后不可能做到“立等可取”吧？}

“立等可取”的意思就是尽量减少顾客的不方便，包括提供一个暂时替代的产品，并不是所有东西都能马上好的意思。（2013-09-07）

\BookQuestion{网友：您当初经营企业时，对售后服务的原则和要求是什么呢？}

我理解的售后其实是销售的延伸，是产品承诺的一部分，企业有义务解决承诺期内产品的所有问题。小霸王学习机的“一年包换”就是基于这个原则定出来的（现在产品周期小于一年，这条可能不再适用了）。（2019-04-27）

% 来源：EPUB part113.xhtml
\subsection{案例}

\BookMinorHeading{技术导向：摩托罗拉}

我觉得摩托罗拉以前是个工程师导向的公司。我最早用的就是摩托罗拉手机，当时特别希望手机上能有个时钟，结果等了好几年都没见加上来，后来就换诺基亚了。我觉得是文化的问题，他们过去不是消费者导向的公司，以后不知能否改过来。（2010-04-05）

% 来源：EPUB part114.xhtml
\BookMinorHeading{利润导向：柯达}

\BookQuestion{“柯达的双重讽刺，日前柯达宣布，为降低成本将逐渐放弃数码相机业务，其实在这一领域，柯达一度遥遥领先，拥有过巨大的优势，但是他们的营销战略从根本上就是错误的，才会走上了今天这条不归路……”（新闻）}

“关键的动机并不是让大家更好地拍照，而是为了卖他们的胶卷。”利润导向和消费者导向的差别长期而言就是生与死的差别。突然想到惠普（HP）打印机有点为了卖墨盒的意思（纯属乱感觉，没经过数据比较）。（2012-04-28）

% 来源：EPUB part115.xhtml
\BookMinorHeading{市占率导向：诺基亚}

\BookQuestion{网友：诺基亚股价从2008年28美元跌到最近6美元，市场占有率维持在三分之一以上，从中给我很多启示。该司的行业领导地位，亮丽的财务报表，不错的全球分销渠道，产品链齐全，很好的净资产收益率等，使其本应是很好的投资标的。但从目前来说，重仓投资者亏大了。该司目前面临iPhone、iPad的高端市场竞争，手机操作系统被苹果颠覆，这些引起管理层的动荡等系列问题。我觉得以过度竞争的企业为投资标的时应非常谨慎，因为这类企业护城河不宽。}

诺基亚太把当前的市场占有率当回事了。

占有率是结果，不是不重要，但不能盲目追求，否则以后会有麻烦。有点像GDP，不是不重要，但不能盲目追求。

\BookQuestion{网友：看来市场占有率和利润一样，只能看成做对事情的结果，如果整天盯着看，动作就会变形。满足消费者需求，创造价值才是正道。诺基亚还犯这样的错，可见要做到真不容易啊。}

说得好。眼光盯住市场占有率确实会导致公司做一些本末倒置的事情。

\BookQuestion{网友：占有率只是个数字，不能阻止OPPO超越诺基亚。}

是否超越谁并不是我们的目标。最重要的是做出好东西给用户。（2010-10-08）

\BookQuestion{网友：段总言下之意是不应该把追求市场占有率放在首位？}

看看iPhone和诺基亚就明白了。（2010-10-10）

\BookQuestion{网友：诺基亚有很多品类产品，低端占市场，中高端挣利润。如果没有苹果，我觉得似乎诺基亚在高端市场还有好日子过。}

这是诺基亚的大错之一。从他们推出低端机开始我就猜他们会有今天了。（2010-10-15）

\BookQuestion{网友：诺基亚是属于做对的事情中做错了事，还是做了错的事情？我觉得诺基亚需要重新回归关注用户体验这条路上。}

诺基亚出的问题大概就是“成功是失败之母”类的。由于这些年一直非常成功，他们大概忘了会成功的原因了。盲目追求市场占有率可能也是原因之一。简单讲，我认为是文化出问题了，也就是没能坚持做对的事情（或者叫做了错的事情）。

诺基亚是一家很有实力的公司，如果现在开始真正反省（或者用你的话就是重新回归到关注用户体验上来），几年内说不定还有卷土重来的机会，再过几年就有点积重难返了。

微软好像也有同样的问题，雅虎其实也是犯了类似的错误。我们旁观者清容易啊，当局者确实很容易犯这类错误。（2011-02-10）

好像很多人太在意市场份额这个其实没太多意义的指标。其实要提高市场份额不难，但保持有意义的份额不容易。诺基亚当年就是太在意市场份额才把自己推到绝境的。（2013-02-24）

\BookQuestion{网友：我不大明白这个“有意义的市场份额”。是不是可以反过来这样理解，无意义的市场份额就如同不能持续产生现金流的账面价值。}

你这样理解也是可以的。（2013-02-26）

\BookQuestion{网友：诺基亚手机为啥死掉了？}

当时诺基亚一年出大概50个机种，岂有不死之理？！

说个小故事，大概2001年前后，我还在读中欧时，曾经跟班上一个诺基亚的同学说过，我觉得目前的诺基亚有蛮大的问题。反映到诺基亚北京总部后，对方说有兴趣听听，让我去。我当时觉得他们的问题为什么要叫我去呢，架子很大嘛，所以我没理会。当然，就算我说了也未必会有用的。当年摩托罗拉的一个中国负责人倒是确实飞过来找我聊过他们的问题，但结论是：你说的有道理，但控制权不在我手里，改变不了啊。

其实无非就是要做对的事情。出50个机种显然是不对的，有时候一个不对的事情就会要了它们的命啊，结果也证实了这点。

\BookQuestion{网友：诺基亚是不是一直没想明白手机是用来上网的，不是用来打电话的？}

诺基亚不是用户导向的，你看到的这个现象其实是他们内在文化的外在表现。（2018-10-23）

\BookQuestion{网友：一般哪些公司会有大麻烦，倒下就起不来那种？}

任何公司倒下去以后都很难再起来，所以倒下再起来的故事才会流传开来，不然就是常态了。（2010-10-10）

居然认为管理层没有重大决策失误。诺基亚是典型的成功是失败之母的典范，失误绝对一箩筐。补充一下：有时候必须决策时不决策也是重大决策失误。（2012-05-19）
